% !TeX root = ../../../金融经济学_王江_讲义及习题.tex
% 教学示意。曲面的五个圆形起伏区域互不重叠，故全局最值恰在最高峰、最低谷的中心。
\begin{center}
\begin{minipage}{\linewidth}
\WJFigureStart{fig:wj-three-extreme-value}%
\centering
\begin{minipage}[t]{.42\linewidth}
\centering
\small （a）一元连续函数\par\smallskip
\begin{tikzpicture}[>=Stealth,font=\small,x=1cm,y=1cm]
 \path[use as bounding box] (-.35,-.65) rectangle (4.65,3.7);
 \coordinate (low) at (.55,.42);
 \coordinate (high) at (2.95,2.68);
 \draw[->,gray!80!black] (0,0)--(4.38,0) node[right] {$x$};
 \draw[->,gray!80!black] (0,0)--(0,3.16) node[above] {$f(x)$};
 \draw[gray,densely dotted] (low)--(.55,0);
 \draw[gray,densely dotted] (3.98,1.55)--(3.98,0);
 \draw[blue!65!black,thick] (low)
  .. controls (.87,.42) and (.83,1.72) .. (1.32,1.72)
  .. controls (1.7,1.72) and (1.73,.94) .. (2.02,.94)
  .. controls (2.40,.94) and (2.45,2.68) .. (high)
  .. controls (3.42,2.68) and (3.54,1.50) .. (3.98,1.55);
 \fill[blue!65!black] (3.98,1.55) circle (1.6pt);
 \fill[red!75!black] (high) circle (2pt) node[above] {全局最大值};
 \fill[red!75!black] (low) circle (2pt) node[right,xshift=2pt] {全局最小值};
 \node[below] at (.55,0) {$a$};
 \node[below] at (3.98,0) {$b$};
\end{tikzpicture}
\end{minipage}\hfill
\begin{minipage}[t]{.55\linewidth}
\centering
\small （b）二元连续函数\par\smallskip
% !TeX root = ../../../金融经济学_王江_讲义及习题.tex
\begin{tikzpicture}[
 >=Stealth,font=\small,
 x={(4.0cm,-.65cm)},y={(-2.05cm,-1.24cm)},z={(0cm,1.34cm)},
 declare function={
  wjextremebump(\dx,\dy,\rad)=max(0,1-((\dx)^2+(\dy)^2)/(\rad)^2)^2;
  wjextremeheight(\u,\v)=.87+1.33*wjextremebump(\u-.56,\v-.44,.21)
   -.68*wjextremebump(\u-.18,\v-.23,.15)+.53*wjextremebump(\u-.12,\v-.68,.14)
   +.49*wjextremebump(\u-.80,\v-.12,.13)-.35*wjextremebump(\u-.44,\v-.12,.12);
 }]
 \filldraw[fill=blue!6,draw=blue!55!black] (0,0,0)--(1,0,0)--(0,1,0)--cycle;
 \draw[->,gray!80!black] (0,0,0)--(1.12,0,0) node[right] {$c_0$};
 \draw[->,gray!80!black] (0,0,0)--(0,1.12,0) node[left] {$c_1$};
 \draw[->,gray!80!black] (0,0,0)--(0,0,2.12) node[above] {$f(c_0,c_1)$};
 % r is distance from the origin in budget shares; s divides that budget between the goods.
 \foreach \i in {0,...,23} {
  \pgfmathsetmacro{\rNear}{\i/24}
  \pgfmathsetmacro{\rFar}{(\i+1)/24}
  \foreach \j in {0,...,23} {
   \pgfmathsetmacro{\sLeft}{\j/24}
   \pgfmathsetmacro{\sRight}{(\j+1)/24}
   \pgfmathsetmacro{\tone}{min(76,14+29*wjextremeheight((\rNear+\rFar)/2*(\sLeft+\sRight)/2,
       (\rNear+\rFar)/2*(1-(\sLeft+\sRight)/2)))}
   \filldraw[fill=blue!\tone!cyan!38!white,draw=blue!46!black,line width=.12pt]
    ({\rNear*\sLeft},{\rNear*(1-\sLeft)},{wjextremeheight(\rNear*\sLeft,\rNear*(1-\sLeft))})--
    ({\rFar*\sLeft},{\rFar*(1-\sLeft)},{wjextremeheight(\rFar*\sLeft,\rFar*(1-\sLeft))})--
    ({\rFar*\sRight},{\rFar*(1-\sRight)},{wjextremeheight(\rFar*\sRight,\rFar*(1-\sRight))})--
    ({\rNear*\sRight},{\rNear*(1-\sRight)},{wjextremeheight(\rNear*\sRight,\rNear*(1-\sRight))})--cycle;
  }
 }
 \coordinate (surfaceMax) at (.56,.44,{wjextremeheight(.56,.44)});
 \coordinate (surfaceMin) at (.18,.23,{wjextremeheight(.18,.23)});
 \draw[red!70!black,densely dashed] (surfaceMax)--(.56,.44,0);
 \fill[red!75!black] (.56,.44,0) circle (1.5pt);
 \fill[red!75!black] (surfaceMax) circle (2pt) node[above right] {全局最大值};
 \fill[red!75!black] (surfaceMin) circle (2pt);
 \draw[red!75!black] (surfaceMin)--([xshift=-13mm,yshift=-7mm]surfaceMin)
   node[below,fill=white,inner sep=1pt] {全局最小值};
 \node[blue!55!black,below right] at (.66,.34,0) {预算边界};
\end{tikzpicture}
%
\end{minipage}
\WJFigureCaption{闭区间和预算三角形上的全局最大值与最小值}
\end{minipage}
\end{center}
